\section{Rotation-Equivariant Conditional Spherical Neural Fields}\label{sec:RECSNF}

We wish to construct a generative model of spherical signals that is rotation equivariant with respect to the latent representation of the signal. That is to say, a rotation of the latent representation corresponds to a rotation of a spherical signal and a signal can be reconstructed with exactly the same accuracy in any rotation. We propose two variants of Vector Neurons \cite{deng_vector_2021} for $SO(2)$ and $SO(3)$ equivariant representation of spherical signals.

\subsection{Spherical Signals as Vector Neurons}
As in Vector Neurons \cite{deng_vector_2021}, we use an ordered list of 3D vectors for our latent representation. In contrast to Vector Neurons \cite{deng_vector_2021}, our signal is defined over directions, not positions. 

Our model takes the form of a conditional spherical neural field, $f:S^2\times\R^{3\times N}\rightarrow \R^M$, such that $f(\mathbf{d},\mathbf{Z})$ computes the value of the signal represented by vector neuron latent code $\mathbf{Z}\in\R^{3\times N}$ in direction $\mathbf{d}\in\mathbb{R}^3$, with $\|\mathbf{d}\|=1$. Our rotation equivariant construction is independent of the specific architecture or conditioning mechanism of $f$. For colour images, $M=3$ and $f$ outputs an RGB colour. By using a spherical neural field, we are agnostic to how the signals are sampled on the sphere. We can generate any sampling simply by choosing the grid of directions as appropriate. We also avoid boundary effects since our domain is continuous. This is a significant advantage over methods for learning spherical signals that operate in the 2D image domain, for example using equirectangular projections.

\subsection{SO(3) Equivariance}

We construct the neural field such that it is \emph{invariant} to a rotation of both $\mathbf{d}$ and $\mathbf{Z}$ simultaneously i.e.
\begin{equation}
    f(\mathbf{R}\mathbf{d},\mathbf{R}\mathbf{Z})=f(\mathbf{d},\mathbf{Z}),\quad \text{with } \mathbf{R}\in SO(3).\label{eqn:invariance}
\end{equation}
This entails that the neural field is \emph{equivariant} with respect to a rotation of $\mathbf{Z}$ only (i.e.~rotating $\mathbf{Z}$ corresponds to rotating the spherical signal such that $f(\mathbf{d},\mathbf{R}\mathbf{Z})=f(\mathbf{R}^\top\mathbf{d},\mathbf{Z})$). See a visualisation of this property on the right of Figure \ref{fig:teaser}.

Key to our approach is a transformation of the inputs to the neural field, $(\mathbf{d},\mathbf{Z})$, such that they are rotation invariant. These divide into two parts:
\begin{enumerate}
    \item $\mathbf{d}^\prime$ - the directional input to the spherical neural field,
    \item $\mathbf{Z}^\prime$ - the latent code on which the neural field is conditioned.
\end{enumerate}

So long as these two inputs satisfy the desired rotation invariance then the neural field itself exhibits this invariance.

The direction in which we wish to evaluate the spherical neural field must be encoded relative to the latent code. 
%in the particular rotation in which we encounter it.
This is satisfied by using the inner product $\langle \mathbf{d},\mathbf{Z} \rangle$, i.e.~the matrix-vector product $\mathbf{d}^\prime=\mathbf{Z}^\top\mathbf{d}\in\R^N$. 
To see that this transformation satisfies the required invariance in \eqref{eqn:invariance}, we note that, for any $\mathbf{R}\in SO(3)$: 
\begin{equation}
(\mathbf{RZ})^\top\mathbf{Rd}=\mathbf{Z}^\top\mathbf{R}^\top\mathbf{Rd}=\mathbf{Z}^\top\mathbf{I}\mathbf{d}=\mathbf{Z}^\top\mathbf{d}.\label{eqn:invar}
\end{equation}
Unlike Vector Neurons, our input is a direction with unit norm, not a position. Hence, the rotation invariant feature $\|\mathbf{d}\|$ conveys no information and we do not use it. The dimensionality of the $\mathbf{d}^\prime$ part of the invariant representation is therefore $N$.

For the latent code, we simply apply the Vector Neurons \cite{deng_vector_2021} invariant layer: $\mathbf{Z}^\prime=\text{VN-In}(\mathbf{Z}) \in \R^{3 \times N}$. Where $\text{VN-In}(\mathbf{Z}_{n}) = \mathbf{Z}_{n}\mathbf{T}_{n}$, with $n$ indexing a single latent vector of $\mathbf{Z}$ and $\mathbf{T}_{n} \in \R^{3 \times 3}$ being an equivariant transform produced from $\mathbf{Z}_{n}$ via a \text{VN-MLP} with an output dimension of $3$. The dimensionality of the $\mathbf{Z}^\prime$ part of the invariant representation is therefore $3N$.

Since our neural field is equivariant \emph{we do not need to augment our training data over the space of rotations}. Observing a spherical signal once in a particular rotation means we can reconstruct it with the same accuracy in any rotation by rotating the latent code.

\subsection{SO(2) Equivariance}\label{sec:SO2}

In certain settings it is desirable to restrict the equivariance to a specific subset of rotations. We propose a restricted transformation of the neural field inputs that are invariant only to rotations, $\mathbf{R}_{\mathbf{a}}\in SO(3)_{\mathbf{a}}\cong SO(2)$, about a given axis $\mathbf{a}\in\mathbb{R}^3$, $\|\mathbf{a}\|=1$, i.e.~a subgroup of $SO(3)$ that is isomorphic to $SO(2)$.

In order to construct the invariant features, we define $\mathbf{b}_{\perp\mathbf{a}}$ as the vector rejection of $\mathbf{b}$ onto $\mathbf{a}$, i.e.~the orthogonal projection of $\mathbf{b}$ onto the plane that is orthogonal to $\mathbf{a}$. Note that this component of $\mathbf{b}$ \emph{is} affected by $\mathbf{a}$-axis rotation. We can write the vector rejection of $\mathbf{b}$ onto $\mathbf{a}$ as a matrix multiplication by first rotating $\mathbf{a}$ onto an arbitrarily chosen standard basis, we choose $\mathbf{e}_x=[1,0,0]^\top$, followed by orthogonal projection and rotation back:
\begin{equation}
    \mathbf{b}_{\perp\mathbf{a}} = \mathbf{R}_{\mathbf{e}_x\rightarrow\mathbf{a}}\text{diag}(0,1,1)\mathbf{R}_{\mathbf{a}\rightarrow\mathbf{e}_x}\mathbf{b},
\end{equation}
where $\text{diag}(0,1,1)$ is the orthogonal projection onto the $y$-$z$ plane and $\mathbf{R}_{\mathbf{a}\rightarrow\mathbf{e}_x}$ is the rotation matrix that rotates $\mathbf{a}$ onto $\mathbf{e}_x$. Since we are not interested in the specific coordinate frame of the invariant features, we simplify by dropping the second rotation and only retaining the $y$ and $z$ coordinates (the $x$ coordinate will be zero in the rotated space). This provides 2D coordinates:
\begin{equation}
    \mathbf{b}_{\perp_\text{2D}\mathbf{a}} = \begin{bmatrix}
 0 & 1 & 0 \\
 0 & 0 & 1
 \end{bmatrix}\mathbf{R}_{\mathbf{a}\rightarrow\mathbf{e}_x}\mathbf{b},
\end{equation}

We also define $\text{proj}_\mathbf{a}\mathbf{b}$ as the scalar projection of $\mathbf{b}$ onto $\mathbf{a}$. This component of $\mathbf{b}$ \emph{is not} affected by $\mathbf{a}$-axis rotation. Again, we can write this in matrix form as:
\begin{equation}
    \text{proj}_\mathbf{a}\mathbf{b} = \left[1,0,0\right]\mathbf{R}_{\mathbf{a}\rightarrow\mathbf{e}_x}\mathbf{b}.
\end{equation}

We can now use these operations to define the $SO(2)$-invariant inputs.

The directional part of the invariant input now contains three components: $\mathbf{d}_\mathbf{a}^\prime=(\text{proj}_\mathbf{a}\mathbf{d}, \langle \mathbf{d}_{\perp_\text{2D}\mathbf{a}},\mathbf{Z}_{\perp_\text{2D}\mathbf{a}} \rangle, \|\mathbf{d}_{\perp_\text{2D}\mathbf{a}}\|)$. The first is the invariant component of $\mathbf{d}$, i.e.~the part that is unaffected by an $\mathbf{a}$-axis rotation and can therefore be used directly. The second encodes $\mathbf{d}$ relative to $\mathbf{Z}$ in the plane perpendicular to $\mathbf{a}$. The third measures the norm of $\mathbf{d}$ projected into the plane perpendicular to $\mathbf{a}$, which is unchanged by rotations about $y$. The dimensionality of this part of the invariant representation is therefore $N+2$.

The conditioning part of the invariant input now contains two components: $\mathbf{Z}_\mathbf{a}^\prime=(\text{proj}_\mathbf{a}\mathbf{Z},\text{VN-In}(\mathbf{Z}_{\perp_\text{2D}\mathbf{a}}))$. The first is simply the invariant component of each column of $\mathbf{Z}$. The second is the Vector Neurons invariant layer transformation of the latent vectors projected into the plane perpendicular to $\mathbf{a}$. The dimensionality of this part of the invariant representation is therefore $3N$.

\subsection{Reflection Equivariance}
\label{sec:reflection_equi}

As we noted in Section \ref{sec:litrevinv}, Vector Neurons are in fact equivariant or invariant to more general $O(3)$ transformations, i.e.~to both rotations \emph{and reflections}. Our $SO(3)$- and $SO(2)$-invariant transformations described above are also invariant to the more general $O(3)$ and $O(3)_{\mathbf{a}}\cong O(2)$ transformation groups respectively. This is because the relationship in \eqref{eqn:invar} also holds for any $\mathbf{R}\in O(3)$.

This means that our $O(3)$ formulation is invariant to reflections about any axis. On the other hand, our $O(2)$ formulation is invariant to reflections about any plane parallel to $\mathbf{a}$ but not to reflections about the plane perpendicular to $\mathbf{a}$. To see this, note that $\text{proj}_\mathbf{a}\mathbf{d}$ would change sign if $\mathbf{d}$ were reflected about the plane perpendicular to $\mathbf{a}$ and therefore $\mathbf{d}^\prime_\mathbf{a}$ would be changed.